% chapters/04_conclusiones_trabajo_futuro.tex
\chapter{Conclusiones y trabajo futuro}
\label{chap:conclusiones_trabajo_futuro}

Este capítulo cierra el trabajo sintetizando las principales conclusiones obtenidas y planteando posibles líneas de continuación. El objetivo no es repetir los resultados del Capítulo~\ref{chap:diseno_experimental}, sino extraer una lectura global sobre la utilidad de los modelos sustitutos, especialmente los procesos gaussianos, para la búsqueda de entradas óptimas en escenarios con evaluaciones limitadas.

\section{Conclusiones}
\label{sec:conclusiones}

El objetivo principal de este TFG ha sido estudiar el uso de procesos gaussianos como modelos sustitutos probabilísticos para aproximar funciones objetivo costosas o desconocidas y utilizarlos como apoyo en la búsqueda de configuraciones prometedoras. Esta elección resulta adecuada porque el GP proporciona tanto una predicción puntual como una estimación de incertidumbre, dos elementos fundamentales cuando no basta con predecir una respuesta, sino que también es necesario decidir dónde conviene invertir nuevas evaluaciones.

La experimentación se ha organizado en dos escenarios complementarios. En los benchmarks sintéticos se ha evaluado el ciclo completo de optimización basada en modelos sustitutos: diseño inicial, ajuste del GP, selección de nuevos puntos mediante Expected Improvement y actualización del conjunto observado. Los resultados muestran que el proceso GP + EI mejora el mejor valor encontrado respecto al diseño inicial en todos los benchmarks analizados, aunque la magnitud de la mejora depende claramente de la dimensión, la dificultad del paisaje objetivo, el kernel utilizado y la información inicial disponible.

Además, los benchmarks han mostrado que la calidad predictiva del modelo y su utilidad para optimizar no son exactamente equivalentes. Algunos modelos pueden obtener un buen MAE sobre el conjunto de test sin ser necesariamente los que mejor guían la búsqueda del mínimo. Por ello, en optimización basada en modelos sustitutos no basta con evaluar únicamente el error medio global: también es necesario analizar cómo la incertidumbre del modelo interactúa con la función de adquisición y si esa combinación permite localizar regiones prometedoras del dominio.

En el caso real de dietas para \textit{Hermetia illucens}, los datos experimentales se han transformado en una representación supervisada basada en variables de formulación, composición nutricional y tipo de subproducto. El protocolo Leave-One-Diet-Out ha permitido evaluar si el modelo generaliza a dietas completas no vistas durante el entrenamiento. Bajo esta validación, el GP compuesto mejora al baseline Dummy en los tres objetivos activos: FCR, quitina y proteína. Esto indica que el modelo captura cierta señal estructural en las variables de formulación y composición de la dieta, y no se limita únicamente a memorizar réplicas de dietas ya observadas.

No obstante, los resultados del caso real deben interpretarse con cautela. La generalización no es uniforme para todas las dietas y la incertidumbre predictiva solo aparece parcialmente calibrada. En quitina, la desviación estándar del GP se relaciona mejor con el error observado, mientras que en FCR y proteína esta relación es más débil. Por tanto, la incertidumbre del GP se considera una señal útil de apoyo, pero no una garantía absoluta de fiabilidad.

Finalmente, el análisis de pre-\emph{infill} permite transformar las predicciones y la incertidumbre del GP en candidatos concretos para una posible evaluación futura. Sin embargo, estos candidatos no deben entenderse como nuevas dietas óptimas demostradas, sino como hipótesis experimentales. Los mejores valores observados siguen correspondiendo a dietas ya ensayadas, mientras que las formulaciones propuestas por EI representan puntos que podrían merecer una nueva evaluación. En conjunto, el trabajo confirma que los modelos sustitutos no sustituyen la validación experimental, pero sí pueden reducir el espacio de búsqueda, cuantificar incertidumbre y priorizar nuevas evaluaciones cuando el presupuesto experimental es limitado.

\section{Trabajo futuro}
\label{sec:trabajo_futuro}

La primera línea de trabajo futuro consiste en cerrar el ciclo experimental del caso real. En este TFG, el caso Entomotive llega hasta una fase de pre-\emph{infill}: el modelo se entrena con los datos disponibles, se valida mediante LODO y se utiliza para proponer candidatos prospectivos. El siguiente paso natural sería seleccionar una o varias formulaciones candidatas, evaluarlas experimentalmente y reentrenar el GP con las nuevas observaciones. De este modo, el caso real pasaría de una validación retrospectiva a un verdadero proceso secuencial de optimización basada en modelos sustitutos.

Una segunda línea de continuación sería formalizar una función de decisión multiobjetivo. En el análisis actual, FCR, quitina y proteína se estudian por separado, lo que permite interpretar cada objetivo de forma individual. Sin embargo, desde un punto de vista productivo, la elección de una dieta depende normalmente de varios criterios simultáneos, como eficiencia de conversión, composición nutricional, coste, disponibilidad del subproducto o viabilidad experimental. Definir una regla de decisión junto con criterio experto permitiría estudiar compromisos entre objetivos y no solo óptimos individuales.

También sería importante ampliar el dataset experimental. El caso real utilizado contiene 11 dietas con 3 réplicas por dieta, suficiente para una primera validación metodológica, pero todavía limitado para extraer conclusiones biológicas fuertes. Incorporar nuevas dietas, más niveles de inclusión y más réplicas permitiría mejorar la estabilidad del modelo, estudiar mejor la calibración de la incertidumbre y comprobar si las relaciones aprendidas se mantienen fuera del rango actualmente observado.

Desde el punto de vista metodológico, futuras extensiones podrían explorar modelos capaces de tratar varios objetivos de forma conjunta, como procesos gaussianos multi-salida o modelos jerárquicos. Esto permitiría aprovechar posibles relaciones entre FCR, proteína y quitina, especialmente en escenarios de pocos datos. También sería interesante estudiar modelos con ruido heterocedástico o técnicas de recalibración probabilística, dado que la incertidumbre del GP no se comporta igual en todos los objetivos.

Por último, el proceso de adquisición podría ampliarse más allá de Expected Improvement. En futuras versiones podrían utilizarse criterios multiobjetivo, criterios con restricciones o estrategias de \emph{batch infill} que propongan varias dietas para ensayar en una misma ronda experimental. Estas extensiones permitirían convertir el marco desarrollado en una herramienta más completa de apoyo a la decisión, capaz de guiar campañas experimentales sucesivas con un uso más eficiente del presupuesto disponible.
